\begin{textsample}{POS Dim 2 – 2019\_c5 – Score -1.00 – 2019\_c5/enade\_2019\_odonto\_2321d6f7bf.txt}  \label{ex:f2_pos_343}
CAPÍTULO 5 DISTRIBUIÇÃO DOS CONCEITOS 5.1 PANORAMA NACIONAL DA DISTRIBUIÇÃO DOS CONCEITOS O Conceito Enade é um indicador de qualidade que avalia os cursos por intermédio dos desempenhos dos estudantes a partir dos resultados do Enade.
Os valores possíveis vão de 1 ( pior situação ) a 5 ( melhor situação ).
Os cursos com apenas um ou sem nenhum concluinte participante não obtêm o Conceito Enade, ficando Sem Conceito ( SC ).
Na Tabela 5.1, são apresentadas a quantidade e a distribuição de cursos de Odontologia participantes do Enade/2019, por faixa de conceito e Grande Região.
Enfatiza -se, mais uma vez, que a diferença entre os cursos tabulados neste capítulo e no capítulo 2 corresponde aos cursos que não foram avaliados, ou seja, àqueles, em princípio, sem alunos concluintes que tivessem sido inscritos para a prova.
Observando -se os dados da Tabela 5.1, nota -se que, dos 238 cursos participantes, 93 ( 39,1 \% ) classificaram -se com conceito 3, o valor modal.
Esse foi também o conceito modal nas regiões Sudeste ( 40,2 \% ), Sul ( 49,0 \% ) e Centro-Oeste ( 45,0 \% ).
Na região Norte a classe modal foi o conceito 2 ( 50,0 \% ), e, na região Nordeste, a classe modal foi o conceito 4 ( 38,2 \% ).
O conceito 4 foi o segundo mais frequente em nível nacional ( 27,7 \%, correspondendo a 66 cursos ), e o conceito 2, o terceiro ( 23,9 \%, correspondendo a 57 cursos ).
Houve, ainda, 16 cursos ( 6,7 \% ) que receberam conceito 5, e cinco cursos ( 2,1 \% ) que receberam conceito 1.
Dos 238 cursos de Odontologia, um ( 0,4 \% ) ficou Sem Conceito ( SC ).
A região Norte participou com 22 cursos ou 9,2 \% do total nacional.
Desses, cinco cursos receberam o conceito 3, o que equivale a 22,7 \% do total regional.
A 11 cursos ( 50,0 \% ) atribuiu -se conceito 2, além de ser o conceito modal na região, como já comentado, e a dois cursos ( 9,1 \% ), conceito 4.
Nenhum curso ficou Sem Conceito ( SC ), 3 cursos receberam conceito 1 ( 13,6 \% ) e um curso recebeu conceito 5 ( 4,5 \% ).
A região Nordeste participou com 55 cursos ou 23,1 \% do total nacional.
Nessa região, 18 cursos ( 32,7 \% em termos regionais ) obtiveram conceito 3.
Os conceitos 2 e 4 foram atribuídos, respectivamente, a 11 e a 21 cursos ( 20,0 \% e 38,2 \%, respectivamente ), tendo o conceito 4 como classe modal, como já comentado.
Ainda nessa região, quatro cursos ( 7,3 \% ) foram avaliados com o conceito 5.
Nenhum curso ficou Sem Conceito ( SC ) e um recebeu conceito 1 ( 1,8 \% ).
Dos 92 cursos participantes da região Sudeste ( 38,7 \% do total nacional ), o conceito 3 foi o conceito modal, como já anteriormente assinalado, tendo sido obtido por 37 cursos ( 40,2 \% ).
O conceito 1 foi atribuído a um curso ( 1,1 \% ), e o conceito 2, a 21 cursos ( 22,8 \% ).
Receberam os conceitos 4 e 5, respectivamente, 25 cursos ( 27,2 \% ) e 8 cursos ( 8,7 \% ).
Nessa região, nenhum curso ficou Sem Conceito ( SC ).
Os 49 cursos da região Sul corresponderam a 20,6 \% do total nacional.
Como já indicado, houve predominância do conceito 3, atribuído a 49,0 \% dos cursos da região ( conceito modal ), ou seja, atribuído a 24 dos 49 cursos participantes na região Sul.
O conceito 2 foi atribuído a nove cursos ( 18,4 \% ) e os conceitos 4 e 5, respectivamente, 13 cursos ( 26,5 \% ) e dois cursos ( 4,1 \% ).
Um curso ficou Sem Conceito ( SC ) na região Sul, e nenhum curso recebeu conceito 1.
Já dos 20 cursos participantes na região Centro-Oeste ( 8,4 \% do total nacional ), nove ( 45,0 \% em termos regionais ) receberam conceito 3, o conceito modal, como já fora destacado.
Cinco cursos ( 25,0 \% ) obtiveram conceito 2, e outros cinco ( 25,0 \% ), conceito 4.
Além disso, nenhum curso ficou Sem Conceito ( SC ).
Nessa região, nenhum curso recebeu o conceito 1 e um curso recebeu o conceito 5 ( 5,0 \% ).
No Gráfico 5.1, é apresentada a distribuição cumulativa do Conceito Enade segundo Grande Região.
Em um gráfico de distribuição cumulativa, em cada valor de abscissa ( eixo x ), o valor apresentado no eixo das ordenadas ( eixo y ) é o valor acumulado da variável ( conceito Enade ).
Observando a linha vertical pontilhada, que ajuda a visualizar o percentual de conceitos acumulados até o valor 3, ou seja, conceitos 1, 2 e 3, pode -se notar que, por exemplo, a região Norte ( linha azul ) apresenta mais de 85 \% dos cursos com conceito 3 ou abaixo deste valor.
A região Nordeste ( linha verde ), por outro lado, apresenta menos de 60 \% dos cursos na mesma situação : com conceito 3 ou abaixo.
Como em todas as regiões houve cursos com conceito 5, as poligonais cumulativas só alcançam 100 \% no conceito máximo.
Em linhas gerais, em um gráfico de distribuição cumulativa, poligonais mais à esquerda representam Grandes Regiões com pior distribuição, e poligonais mais à direita, Grandes Regiões com melhores conceitos.
A região Norte ( linha azul ) apresenta os cursos com conceitos mais baixos, ao passo que as regiões Nordeste ( linha verde ) e Sudeste ( linha vermelha ) apresentam as distribuições com valores maiores.
Cursos em IES situadas nas regiões Sul ( linha rosa ) e Centro-Oeste ( linha laranja ) apresentam uma situação intermediária. 5.2 CONCEITOS POR CATEGORIA ADMINISTRATIVA E GRANDE REGIÃO Na Tabela 5.2, é apresentada a distribuição dos cursos de Odontologia participantes do Enade/2019 por Categoria Administrativa, de acordo com os conceitos por eles alcançados, segundo as Grandes Regiões brasileiras.
Dos 238 cursos participantes, 60 ( 25,2 \% ) eram ministrados em Instituições Públicas, e 178 ( 74,8 \% ), em Instituições Privadas.
De acordo com as informações da Tabela 5.2, em termos nacionais, entre Instituições Públicas e Privadas, observa -se que, dos 16 cursos avaliados com conceito 5, 12 eram oferecidos em IES Públicas, e quatro, em IES Privadas.
Dos 60 cursos participantes de IES Públicas, o conceito 4 foi o valor modal, atribuído a 38 cursos ( 63,3 \% ).
Dos demais cursos de Instituições Públicas participantes, três cursos foram avaliados com conceito 2, cinco cursos ( 8,3 \% ) foram avaliados com conceito 3 e, como já comentado, 12 cursos ( 20,0 \% ), com conceito 5.
Nenhum curso ficou Sem Conceito ( SC ).
Nesta categoria, dois cursos receberam conceito 1. conceito 5.
Nenhum curso ficou Sem Conceito ( SC ).
Nesta categoria, dois cursos receberam conceito 1.
Na rede privada, o conceito modal foi 3, atribuído a 88 cursos, o correspondente a 49,4 \% dos 178 cursos da categoria.
Dos demais cursos participantes, três ( 1,7 \% ) receberam conceito 1, e 54 ( 30,3 \% ), conceito 2.
O conceito 4 foi atribuído a 28 cursos ( 15,7 \% ), e o conceito 5, a 16 cursos ( 2,2 \% ).
Nesta Categoria Administrativa, um curso ( 0,6 \% ) ficou Sem Conceito ( SC ).
Dos 238 cursos participantes, o conceito 3 foi o valor modal, atribuído a 93 cursos ( 39,1 \% ).
Dos demais cursos, cinco cursos ( 2,1 \% ) receberam conceito 1, 57 cursos ( 23,9 \% ) receberam conceito 2, 66 cursos ( 27,7 \% ), conceito 4, 16 cursos ( 6,7 \% ), conceito 5 e, um curso ( 0,4 \% ) ficou Sem Conceito ( SC ).
Na análise por região, observa -se que, na região Norte, as Instituições Privadas participaram com 18 cursos ( 81,8 \% do total regional ), dos quais nenhum ficou Sem Conceito ( SC ), e nenhum recebeu os conceitos 4 e 5.
O conceito modal para as Instituições Privadas na região foi 2, atribuído a 11 cursos, correspondendo a 61,1 \% dessa categoria na região.
Dois cursos foram avaliados com conceito 1, e cinco cursos, com conceito 3.
As Instituições Públicas participaram com quatro cursos na região Norte ( 18,2 \% do total regional ), aos quais foram atribuídos os conceitos 1, 4 e 5.
Nenhum curso ficou Sem Conceito ( SC ) ou com conceito 2 e 3 nessa combinação de categoria e região.
Na região Nordeste, a rede privada concentrou 35 dos 55 cursos participantes, o equivalente a 63,6 \% do total da região.
Dentre os cursos oferecidos por IES Privadas no Nordeste, 17 cursos foram avaliados com conceito 3, o conceito modal.
Receberam conceito 2, 11 cursos, e outros seis cursos, conceito 4.
Nessa combinação de categoria e região, nenhum curso ficou Sem Conceito ( SC ) e nenhum recebeu conceito 5.
As Instituições Públicas da região Nordeste participaram com 20 cursos ( 36,4 \% do total da região ).
Desses, 15 obtiveram conceito 4, o valor modal, seguidos por outros quatro cursos que receberam conceito 5 e um curso que recebeu o conceito 3.
Nenhum curso ficou Sem Conceito ( SC ) ou com os demais conceitos nessa região.
Na região Sudeste, a proporção de cursos da rede privada, 78,3 \%, correspondendo a 72 dos 92 cursos participantes.
Nessa categoria e região, o conceito modal foi 3, atribuído a 35 cursos.
Os demais foram avaliados com conceito 2 ( 19 cursos ), conceito 4 ( 14 cursos ) e conceito 5 ( quatro cursos ), nenhum curso ficou Sem Conceito ( SC ) ou com conceito 1.
Dos 20 cursos oferecidos por Instituições Públicas na região Sudeste, o conceito modal foi 4, atribuído a 11 cursos.
Nessa região, um curso recebeu conceito 1, dois cursos receberam conceito 2, dois cursos receberam conceito 3 e quatro cursos receberam conceito 4.
Nessa região, nenhum curso ficou Sem Conceito ( SC ).
As Instituições Privadas concentraram 37 dos 49 cursos participantes da região Sul, 75,5 \% do total regional.
Desses, 22 cursos obtiveram conceito 3, o conceito modal.
Nesta combinação de Categoria Administrativa e Grande Região, nove cursos receberam conceito 2, outros cinco, conceito 4.
Nenhum curso recebeu conceito 1 ou 5, e um curso ficou Sem Conceito ( SC ).
As Instituições Públicas na região Sul participaram com 12 cursos ( 24,5 \% dos cursos da região ), aos quais foram atribuídos os conceitos 3, 4 ou 5.
Nenhum curso da região Sul oferecido por IES Pública ficou Sem Conceito ( SC ) ou recebeu os conceitos 1 e 2.
Na região Centro-Oeste, 16 dos 20 cursos participantes eram de Instituições Privadas ( 80,0 \% em termos regionais ).
Desses, nove receberam conceito 3, o conceito modal.
Dos demais cursos, quatro receberam conceito 2, e outros três, conceito 4.
Nenhum curso ficou Sem Conceito ( SC ), e nenhum curso recebeu conceito 1 ou 5.
Quanto aos quatro cursos oferecidos por Instituições Públicas na região Centro-Oeste ( 20,0 \% do total regional ), dois cursos apresentaram conceito 4, conceito modal, um recebeu conceito 2 e um recebeu conceito 5.
Nenhum curso ficou Sem Conceito ( SC ) ou recebeu os demais conceitos nessa combinação de categoria e região.
No Gráfico 5.2, apresenta -se a distribuição do Conceito Enade, segundo a Categoria Administrativa da IES.
Os cursos em IES Públicas ( linha azul ) apresentam uma poligonal mais à direita do que os em IES Privadas ( linha verde ) e, consequentemente, uma distribuição de Conceitos Enade com valores maiores. 5.3 CONCEITOS POR ORGANIZAÇÃO ACADÊMICA E GRANDE REGIÃO Na Tabela 5.3, encontra -se a distribuição dos conceitos atribuídos aos cursos participantes do Enade/2019, na Área de Odontologia, por Organização Acadêmica, segundo as Grandes Regiões brasileiras.
Dos 238 cursos de Odontologia participantes, 136 eram oferecidos em Universidades, 57 em Centros Universitários e 45 em Faculdades, nenhum curso foi oferecido em CEFET/IF.
Esta distribuição corresponde a, respectivamente, 57,1 \%, 23,9 \%, 18,9 \% e 0,0 \% dos cursos.
De acordo com os dados apresentados, de todos os 16 cursos avaliados com conceito 5, 15 eram vinculados a Universidades.
Os cursos vinculados a esse tipo de Organização Acadêmica tiveram o conceito 4 como conceito modal, atribuído a 51 cursos ( 37,5 \% ).
Os demais cursos vinculados a Universidades receberam os conceitos 1 ( dois cursos ), 2 ( 22 cursos ), 3 ( 45 cursos ) e conceito 5 ( 15 cursos, como já mencionado ).
Ficou Sem Conceito ( SC ), um curso.
Dos cursos em Centros Universitários, o conceito modal foi 3, atribuído a 34 cursos ( 59,6 \% ).
Nenhum curso vinculado a esse tipo de Organização Acadêmica ficou Sem Conceito ( SC ), e os outros cursos receberam os conceitos 2 ( 13 cursos ), 4 ( dez cursos ) e nenhum curso recebeu conceito 1 ou 5.
Dos 45 cursos mantidos por Faculdades, 22 ( 48,9 \% ) receberam conceito 2, o conceito modal.
O segundo conceito mais frequentemente atribuído aos cursos oferecidos por essa categoria de IES foi o conceito 3, recebido por 14 cursos.
Na sequência, cinco cursos obtiveram conceito 4, três cursos, conceito 1 e um curso, conceito 5.
Dos cursos oferecidos em Faculdades, nenhum ficou Sem Conceito ( SC ).
Considerando -se, separadamente, as regiões brasileiras, verifica -se que, na região Norte, as Universidades participaram com apenas seis dos 22 cursos avaliados.
Os conceitos 1, 2, 3 e 5 foram atribuídos a um curso cada, e o conceito 4, dois cursos.
Nenhum curso ficou Sem Conceito ( SC ).
Os Centros Universitários da região Norte foram representados por seis cursos, os quais obtiveram conceitos 2 ou 3.
Nenhum curso ficou Sem Conceito ( SC ) ou recebeu os demais conceitos nesta combinação de categoria e região.
As Faculdades participaram com Na região Norte, as Universidades participaram com apenas seis dos 22 cursos avaliados.
Os conceitos 1, 2, 3 e 5 foram atribuídos a um curso cada, e o conceito 4, dois cursos.
Nenhum curso ficou Sem Conceito ( SC ).
Os Centros Universitários da região Norte foram representados por seis cursos, os quais obtiveram conceitos 2 ou 3.
Nenhum curso ficou Sem Conceito ( SC ) ou recebeu os demais conceitos nesta combinação de categoria e região.
As Faculdades participaram com dez cursos na região Norte, esses receberam os conceitos 1 ( dois cursos ) e 2 ( oito cursos, conceito modal ).
Nenhum curso oferecido em Faculdades da região Norte ficou Sem Conceito ( SC ) ou recebeu os conceitos 1, 4 e 5.
Na região Nordeste, as Universidades participaram com 24 dos 55 cursos da Área de Odontologia oferecidos.
O conceito modal 4 foi atribuído a 16 cursos.
Os oito cursos restantes obtiveram conceitos 3 e 5.
Nenhum curso oferecido por Universidades no Nordeste recebeu conceito 1 e 2, e nenhum curso ficou Sem Conceito ( SC ).
Os Centros Universitários contaram com 19 cursos participantes na região Nordeste, nove dos quais receberam o conceito modal 3.
Para os demais cursos, foram atribuídos os conceitos 2 ( seis cursos ) e 4 ( quatro cursos ).
Nenhum curso oferecido por Centros Universitários no Nordeste ficou Sem Conceito ( SC ) ou recebeu os conceitos 1 e 5.
As Faculdades foram representadas por 12 cursos na região Nordeste, dos quais cinco receberam conceito 2 e outros cinco receberam conceito 3, sendo esses os conceitos modais.
Cada um dos conceitos 1 e 4 foi atribuído a um curso.
Na região Nordeste, nenhum curso mantido por Faculdades ficou Sem Conceito ( SC ) ou recebeu conceito 5.
Na região Sudeste, as Universidades concentraram 61 dos 92 cursos de Odontologia da região.
Dos cursos oferecidos em Universidades, nessa região, o conceito modal foi 3, atribuído a 20 cursos.
Os demais cursos receberam os conceitos 1 ( um curso ), 2 ( 16 cursos ), 4 ( 17 cursos ) e 5 ( sete cursos ).
Nenhum curso ficou Sem Conceito ( SC ).
Os Centros Universitários participaram com 18 cursos na região Sudeste, dos quais 11 obtiveram o conceito modal 3.
O conceito 2 foi atribuído a três cursos e o conceito 4, a quatro cursos.
Nenhum curso ficou Sem Conceito ( SC ) ou recebeu conceitos 1 e 5 nessa combinação de organização acadêmica e região.
Já as Faculdades foram representadas por 13 cursos na região Sudeste.
O conceito modal 3 foi atribuído a seis cursos.
Aos demais cursos foram atribuídos os conceitos 2 ( dois cursos ), 4 ( quatro cursos ) e 5 ( um curso ).
Nenhum curso ficou Sem Conceito ( SC ) ou recebeu conceito 1.
Dos 49 cursos da região Sul, 36 eram mantidos por Universidades, dos quais o conceito modal foi 3, atribuído a 17 cursos.
Os demais cursos receberam os conceitos 2 ( três cursos ), 4 ( 13 cursos ) e 5 ( dois cursos ).
Nesse tipo de organização, um curso da região Sul ficou Sem Conceito ( SC ) e nenhum recebeu conceito 1.
Dos seis cursos participantes de Centros Universitários da região Sul, nenhum ficou Sem Conceito ( SC ), e nenhum recebeu os conceitos 1, 4 e 5.
Os conceitos restantes, 2 e 3, foram atribuídos, respectivamente, a um curso e cinco cursos.
Quanto aos sete cursos vinculados a Faculdades na região Sul, cinco receberam o conceito modal 2.
O conceito 3 foi atribuído a dois cursos.
Nenhum curso ficou Sem Conceito ( SC ) ou recebeu conceito 1, 4 e 5.
Na região Centro-Oeste, nove dos 20 cursos oferecidos eram mantidos por Universidades.
Para este tipo de Organização Acadêmica, os conceitos 3 e 4 foram atribuídos a três cursos cada, dois cursos receberam conceito 2 e um curso recebeu conceito 5.
Nesse tipo de organização, nenhum curso ficou Sem Conceito ( SC ) ou recebeu conceito 1.
Os Centros Universitários da região Centro-Oeste contaram com oito cursos, cinco avaliados com conceito 3, classe modal.
Um curso recebeu conceito 2 e dois cursos receberam conceito 4.
Nenhum curso ficou Sem Conceito ( SC ) ou recebeu conceito 1 ou 5.
Dos três cursos oferecidos em Faculdades na região Centro-Oeste, dois receberam o conceito modal 2, e um obteve conceito 3.
Nesta região, nenhum curso ficou Sem Conceito ( SC ), e nenhum obteve conceito 1, 4 e 5.
No Gráfico 5.3, apresenta -se a distribuição do Conceito Enade dos cursos de Odontologia, segundo a Organização Acadêmica da IES.
Os cursos em Universidades ( linha azul ) apresentam uma poligonal mais à direita do que as demais, denotando uma melhor distribuição de Conceitos Enade.
No outro extremo, com distribuição muito semelhante entre si, encontram -se os cursos mantidos por Centros Universitários ( linha verde ) e Faculdades ( linha vermelha ), apresentando as respectivas poligonais mais à esquerda e, consequentemente, distribuição de Conceitos Enade com menores valores.
Destaca -se que a poligonal dos Centros Universitários ( linha verde ) alcança 100 \% no conceito 4 e isso ocorre porque nenhum curso dessa Organização Acadêmica recebeu conceito 5.

% matched lemmas: 
\end{textsample}
